\chapter{Sigma-Point Filters}
\label{ch:bf-sigma-point-filters}

Sigma-point filters replace local linearization with deterministic quadrature.
They propagate a finite set of points through nonlinear maps and reconstruct
approximate Gaussian moments from weighted averages.  This makes them more
accurate than the EKF for some nonlinear transformations, while preserving the
Gaussian filtering closure used by the prediction-error likelihood.  The price
is that the likelihood remains approximate and the derivative path becomes
more complicated.

\section{Moment-Matching Interface}
\label{sec:bf-sigma-point-interface}

At each time step, a sigma-point backend starts from a Gaussian approximation
\[
  x_t \mid y_{1:t-1}
  \approx \calN(m_{t|t-1},P_{t|t-1}).
\]
A deterministic rule constructs points
$\chi^{(j)}_{t|t-1}$ and weights $w_j^{(m)},w_j^{(c)}$.  Passing the points
through the observation map gives
\[
  z_t^{(j)}=h_\theta(\chi^{(j)}_{t|t-1}).
\]
The backend then forms approximate predicted observations, innovation
covariances, and cross covariances:
\begin{align}
  \bar z_t &= \sum_j w_j^{(m)}z_t^{(j)}, \\
  P_{zz,t} &= \sum_j w_j^{(c)}
    (z_t^{(j)}-\bar z_t)(z_t^{(j)}-\bar z_t)^\top + R_\theta, \\
  P_{xz,t} &= \sum_j w_j^{(c)}
    (\chi^{(j)}_{t|t-1}-m_{t|t-1})
    (z_t^{(j)}-\bar z_t)^\top .
\end{align}
The gain is $K_t=P_{xz,t}P_{zz,t}^{-1}$ and the innovation is
$v_t=y_t-\bar z_t$.  The same Gaussian innovation likelihood formula is then
used, with $P_{zz,t}$ replacing $S_t$.

\section{Conjugate Unscented Transform Rules}
\label{sec:bf-cut-rules}

The ordinary unscented transform is not the only deterministic quadrature rule
that can be used in the interface above.  The conjugate unscented transform
(CUT) of \citet{adurthi2012cut,adurthi2012cutfiltering} constructs higher-order
Gaussian point sets by combining symmetric conjugate point families.  The
experimental file
\[
  \texttt{/home/chakwong/python/src/dsge\_hmc/filters/CUTSRUKF.py}
\]
implements a deprecated TensorFlow square-root CUT4 demonstration for a
coordinated-turn tracking model.  It is not wired into the current BayesFilter
filter hierarchy, but it is a useful design reference because it separates the
sigma-point rule from the square-root covariance backend.

For a standard normal variable $Z\in\mathbb{R}^n$, the CUT4-G rule used there
has
\[
  N_{\mathrm{CUT4}} = 2n + 2^n
\]
points: $2n$ axis points and $2^n$ hypercube-corner points.  In the convention
used by the demo, for $n\ge3$,
\[
  r_a^2=\frac{n+2}{2},\qquad
  r_c^2=\frac{n+2}{n-2},
\]
with weights
\[
  w_a=\frac{4}{(n+2)^2},\qquad
  w_c=\frac{(n-2)^2}{2^n(n+2)^2}.
\]
More explicitly, define
\[
  \mathcal A=\{\pm r_a e_i:i=1,\ldots,n\},
  \qquad
  \mathcal C=\{r_c s:s\in\{-1,+1\}^n\}.
\]
The CUT4-G quadrature operator is
\begin{equation}
\label{eq:bf-cut4-quadrature}
  \mathcal Q_{\mathrm{CUT4}}[\varphi]
  =
  w_a\sum_{\xi\in\mathcal A}\varphi(\xi)
  +w_c\sum_{\xi\in\mathcal C}\varphi(\xi).
\end{equation}

\begin{proposition}[CUT4-G Gaussian moment exactness]
\label{prop:bf-cut4-moment-exactness}
For $n\ge3$, the point set
\eqref{eq:bf-cut4-quadrature} integrates every standard-normal monomial in
$Z\in\mathbb R^n$ of total degree at most five exactly.  Equivalently, it is a
degree-five Gaussian rule.
\end{proposition}

\begin{proof}
Both $\mathcal A$ and $\mathcal C$ are centrally symmetric and sign-symmetric
in every coordinate.  Hence every monomial with at least one odd coordinate
exponent has zero quadrature value, matching the standard normal.  This covers
all odd total degrees through five and also the mixed even-odd monomials.

It remains to check total degrees $0$, $2$, and $4$.  The total mass is
\[
  2n w_a+2^n w_c=1 .
\]
For a coordinate $i$,
\[
  \mathcal Q_{\mathrm{CUT4}}[z_i^2]
  =
  2w_a r_a^2+2^n w_c r_c^2
  =
  1
  =
  \mathbb E[Z_i^2],
\]
and, for $i\ne j$,
$\mathcal Q_{\mathrm{CUT4}}[z_i z_j]=0=\mathbb E[Z_iZ_j]$ by sign symmetry.  The
fourth marginal moment is
\[
  \mathcal Q_{\mathrm{CUT4}}[z_i^4]
  =
  2w_a r_a^4+2^n w_c r_c^4
  =
  3
  =
  \mathbb E[Z_i^4],
\]
while the mixed fourth moment is
\[
  \mathcal Q_{\mathrm{CUT4}}[z_i^2z_j^2]
  =
  2^n w_c r_c^4
  =
  1
  =
  \mathbb E[Z_i^2Z_j^2],
  \qquad i\ne j.
\]
All remaining fourth-degree monomials contain an odd coordinate exponent and
therefore vanish on both sides.  These identities exhaust the Gaussian
monomials of total degree at most five, so the rule has degree-five exactness.
\end{proof}

Thus the rule is symmetric, has positive weights, and integrates Gaussian
polynomials through degree five.  By contrast, the usual $2n+1$ unscented rule
is commonly described as third-degree accurate for Gaussian moment propagation.
The extra corner points are what allow CUT4-G to match cross moments such as
$\mathbb{E}[Z_i^2Z_j^2]$ that axis-only rules do not represent in the same way.

The price is exponential growth in the corner set.  CUT4-G uses $42$ points at
$n=5$, $1044$ points at $n=10$, and becomes unattractive in large structural
states unless the quadrature is applied to a low-dimensional stochastic block.
This reinforces the structural lesson of
Chapter~\ref{ch:bf-structural-deterministic-dynamics}: higher-order quadrature
should be spent on the declared stochastic integration coordinates, not on
deterministic-completion directions that add no stochastic law.

\section{Approximation Boundary}
\label{sec:bf-sigma-point-approximation-boundary}

BayesFilter should label sigma-point filters as approximate likelihood
backends unless a special model structure supplies an exactness argument.
The source SR-UKF material contains model-specific claims about quadratic
observation equations and pruned DSGE structure.  Those claims are useful
case-study material, but they should not be promoted to a general BayesFilter
theorem without a separate derivation in BayesFilter notation and a regression
test against an exact reference.

For HMC, this distinction matters.  If the sigma-point likelihood is used in
the target, the chain targets the sigma-point posterior.  If the sigma-point
filter is only a surrogate for an exact nonlinear likelihood, then a correction
or delayed-acceptance policy is needed before exact-posterior claims are made.

\section{Derivative and Compilation Requirements}
\label{sec:bf-sigma-point-derivative-requirements}

Sigma-point HMC requires differentiating through:
\begin{itemize}
  \item the covariance factor used to generate points;
  \item the nonlinear transition and observation maps;
  \item the weighted covariance assembly;
  \item linear solves or factorizations for $P_{zz,t}$;
  \item any PSD projection, jitter, clipping, or fallback branch.
\end{itemize}
The derivative provider must return the derivative of the same approximate log
likelihood that supplies the Metropolis value.  Static shapes are also
non-negotiable: the number of sigma points, state dimension, observed-data
mask shape, and scan length must not vary inside the compiled transition.

\section{Evidence Ladder}
\label{sec:bf-sigma-point-evidence-ladder}

The minimum evidence ladder for a sigma-point backend is:
\begin{enumerate}[label=(\roman*)]
  \item exact recovery on affine Gaussian systems against the Kalman backend;
  \item finite value and finite gradient on controlled nonlinear systems;
  \item eager/compiled parity where a compiled sampler is claimed;
  \item comparison to a denser quadrature, long simulation, or trusted
    reference on low-dimensional nonlinear examples;
  \item HMC diagnostics that remain explicitly separated from convergence
    proof.
\end{enumerate}
This ladder is the consolidation of the source-project SVD/HMC validation
plans, not a completed BayesFilter certification result.

CUT backends should pass the same ladder plus one rule-specific check: for
standard-normal polynomial test functions up to the declared degree, the
implemented points and weights should reproduce the analytic Gaussian moments
before the points are used in a filter.  This catches mistakes in point
enumeration, weights, dimension restrictions, and transformed-coordinate
scaling before nonlinear filtering errors obscure the quadrature issue.
